\documentclass[10pt,a4paper]{amsart}
\usepackage[margin=19mm]{geometry}
\usepackage{amsmath,amssymb,mathtools}
\usepackage[scr=rsfs,cal=euler]{mathalfa}
\usepackage{xfrac}
\usepackage[unicode,hidelinks,bookmarks=false]{hyperref}
\newcommand{\ie}{i.e.\ }
\newcommand{\cf}{cf.\ }
\newcommand{\resp}{respectively\ }
\newcommand{\Subsection}{\subsection}
\providecommand{\nobreakdash}{\nobreak\hbox{-}}
\setlength{\emergencystretch}{2em}
\multlinegap0pt
\usepackage{bm}
\usepackage{mathtools}
\newtheorem{theorem}{Theorem}
\title{Geometric multiscale analysis via nonstationary subdivision schemes}
\author{Hadar Landau, Wael Mattar, Nir Sharon}
\date{Source: arXiv:2507.09668v2 (CC BY 4.0)}
\begin{document}
\maketitle

\noindent\emph{Source and licence.} Geometric multiscale analysis via nonstationary subdivision schemes. Version arXiv:2507.09668v2 (CC BY 4.0), distributed by arXiv under the Creative Commons Attribution 4.0 International licence, http://creativecommons.org/licenses/by/4.0/, as recorded in the arXiv licence metadata for this exact version and shown on the abstract page https://arxiv.org/abs/2507.09668v2. Full licence notice: \texttt{licenses/arxiv-2507.09668-CC-BY-4.0.txt}.\par\smallskip

\noindent\textit{Editorial scope.} Counted unit: the article's theorem decay (source0.tex lines 272--278). The article's complete original proof of it is delivered with it (source0.tex lines 280--315, one \texttt{proof} environment of 36 lines). No mathematics is written, repaired or re-derived: the statement and the proof are the article's own lines, in the article's own notation, and no retained line is substituted or re-flowed.  Retained uncounted as the source-local context the proof needs: lines 208--210.  Nothing was omitted from any retained span, so the package carries no omission record.  No retained line contains a citation, so the wrapper carries no bibliography and no bibliography entry.  No retained line contains a figure environment, an image-inclusion command or drawing code, so no image is delivered and the package declares no asset dependency.  Editorial formatting only, in addition: an AMS \texttt{amsart} preamble that restates the article's own declarations and macros used by the retained lines, namely the environments \texttt{theorem} on separate counters, together with the standard packages those lines need.  The retained environments print in this document as Theorem 1 in order of appearance, whereas the article prints the same items with the numbering of its own full document.  The complete native source of the article as distributed in the arXiv e-print for this exact version is bundled unchanged as \texttt{sources/181499/source0.tex}.
\begin{equation} \label{NS multiscale transform}
    \boldsymbol{c}^{(\ell-1)}=\mathcal{D}_{\boldsymbol{\gamma}^{(\ell)}}\boldsymbol{c}^{(\ell)}, \quad \boldsymbol{d}^{(\ell)}=\boldsymbol{c}^{(\ell)}-\mathcal{S}_{\boldsymbol{\alpha}^{(\ell)}}\boldsymbol{c}^{(\ell-1)}, \quad \ell=1,\dots,J,
\end{equation}

\begin{theorem}[Decay of the detail coefficients]\label{thm. details decay}
Let $f\colon \mathbb{R} \rightarrow \mathbb{R}$ be a bounded, continuously differentiable function with a bounded derivative, and let $\boldsymbol{c}^{(J)}$, $J \in \mathbb{N}$, be its samples over the dyadic grid $2^{-J}\mathbb{Z}$, that is, $\boldsymbol{c}^{(J)}=f|_{2^{-J}\mathbb{Z}}$. Let $\{\mathcal{S}_{\boldsymbol{\alpha}^{(\ell)}} \mid \ell=1,\dots,J\}$ be a family of nonstationary subdivision schemes satisfying $\sum_{i}\alpha^{(\ell)}_{2i}=\sum_{i}\alpha^{(\ell)}_{2i+1}=1$, and let $\{\mathcal{D}_{\boldsymbol{\gamma}^{(\ell)}} \mid \ell=1,\dots,J\}$ be its corresponding family of decimation operators, associated with sequences $\boldsymbol{\gamma}^{(\ell)}$ satisfying $\sum_i \gamma^{(\ell)}_i = 1$ and having finite first moments, that is, $\sum_{i\in\mathbb{Z}}|i|\,|\gamma^{(\ell)}_i|<\infty$. Then the detail coefficients obtained by the multiscale transform~\eqref{NS multiscale transform} satisfy
\begin{equation}\label{details bound}
    \|\boldsymbol{d}^{(\ell)}\|_\infty \leq \bigg(K_{\boldsymbol{\alpha}^{(\ell)},\boldsymbol{\gamma}^{(\ell)}} \cdot \|f^\prime\|_\infty \cdot \prod_{m=\ell+1}^J \|\boldsymbol{\gamma}^{(m)}\|_1 \bigg) \; 2^{-\ell}, \quad \ell=1,\dots,J,
\end{equation}
where $K_{\boldsymbol{\alpha}^{(\ell)},\boldsymbol{\gamma}^{(\ell)}}=K_{\boldsymbol{\gamma}^{(\ell)}}\|\boldsymbol{\alpha}^{(\ell)}\|_1+ K_{\boldsymbol{\alpha}^{(\ell)}}\|\boldsymbol{\gamma}^{(\ell)}\|_1$, with $K_{\boldsymbol{\gamma}^{(\ell)}}=2\sum_{i\in\mathbb{Z}} |\gamma^{(\ell)}_i|\,|i|$ and $K_{\boldsymbol{\alpha}^{(\ell)}}=\sum_{i\in\mathbb{Z}} |\alpha^{(\ell)}_i|\,|i|$, and with the convention that the empty product, occurring for $\ell=J$, equals $1$.
\end{theorem}

\begin{proof}
    Throughout, for a bounded sequence $\boldsymbol{c}$ we write $\Delta\boldsymbol{c}\mathrel{:=}\sup_{k\in\mathbb{Z}}|c_{k+1}-c_k|$ for its maximal first difference; note that $\Delta$ is subadditive and satisfies $\Delta(\boldsymbol{\gamma}*\boldsymbol{c})\leq\|\boldsymbol{\gamma}\|_1\,\Delta\boldsymbol{c}$ for any absolutely summable sequence $\boldsymbol{\gamma}$, since the difference operator commutes with convolution. We first show that $\|\boldsymbol{d}^{(\ell)}\|_\infty$ is proportional to $\Delta\boldsymbol{c}^{(\ell)}$ for all levels $\ell=1,\dots,J$, namely
    \begin{equation}\label{eqn:details_Delta}
        \|\boldsymbol{d}^{(\ell)}\|_\infty \leq K_{\boldsymbol{\alpha}^{(\ell)},\boldsymbol{\gamma}^{(\ell)}} \Delta \boldsymbol{c}^{(\ell)}.
    \end{equation}
    To this end, we begin by evaluating a general term in $\boldsymbol{d}^{(\ell)}$. Observe
    \begin{equation*}
    \begin{split}
        d^{(\ell)}_k & =c^{(\ell)}_k-\sum_i \alpha^{(\ell)}_{k-2i}c^{(\ell-1)}_i=\sum_i \alpha^{(\ell)}_{k-2i}(c^{(\ell)}_k-c^{(\ell-1)}_i)\\
        & =\sum_i \alpha^{(\ell)}_{k-2i}\bigg(c^{(\ell)}_k-\sum_n \gamma^{(\ell)}_{i-n} c^{(\ell)}_{2n}\bigg)\\
        & =\sum_i \alpha^{(\ell)}_{k-2i}\bigg(\sum_n \gamma^{(\ell)}_{i-n}c^{(\ell)}_k-\sum_n \gamma^{(\ell)}_{i-n} c^{(\ell)}_{2n}\bigg)\\
        & =\sum_i \alpha^{(\ell)}_{k-2i}\bigg(\sum_n \gamma^{(\ell)}_{i-n}(c^{(\ell)}_k-c^{(\ell)}_{2n})\bigg).
    \end{split}
    \end{equation*}
    Consequently,
    \begin{equation*}
    \begin{split}
        |d^{(\ell)}_k| & \leq \sum_i|\alpha^{(\ell)}_{k-2i}|\bigg(\sum_n |\gamma^{(\ell)}_{i-n}|\cdot |c^{(\ell)}_k-c^{(\ell)}_{2n}|\bigg)\\ & \leq \sum_i|\alpha^{(\ell)}_{k-2i}|\bigg(\sum_n |\gamma^{(\ell)}_{i-n}|\cdot |2n-k| \cdot \Delta\boldsymbol{c}^{(\ell)}\bigg)\\
        & \leq \sum_i|\alpha^{(\ell)}_{k-2i}|\bigg(\sum_n |\gamma^{(\ell)}_{i-n}|\cdot \big(|2n-2i|+|2i-k|\big) \bigg) \cdot \Delta\boldsymbol{c}^{(\ell)}\\
        & \leq \sum_i|\alpha^{(\ell)}_{k-2i}|\bigg(K_{\boldsymbol{\gamma}^{(\ell)}}+|2i-k| \|\boldsymbol{\gamma}^{(\ell)}\|_1\bigg) \cdot \Delta\boldsymbol{c}^{(\ell)}\\
        & \leq \Big(K_{\boldsymbol{\gamma}^{(\ell)}}\|\boldsymbol{\alpha}^{(\ell)}\|_1+K_{\boldsymbol{\alpha}^{(\ell)}} \|\boldsymbol{\gamma}^{(\ell)}\|_1\Big) \cdot \Delta\boldsymbol{c}^{(\ell)},\\
    \end{split}
    \end{equation*}
    where the last inequality follows from $\sum_i|\alpha^{(\ell)}_{k-2i}|\leq\|\boldsymbol{\alpha}^{(\ell)}\|_1$ and $\sum_i|\alpha^{(\ell)}_{k-2i}|\,|k-2i|\leq K_{\boldsymbol{\alpha}^{(\ell)}}$, since for a fixed $k$ the summation runs over a single parity class of the mask.

    Applying the supremum over $k\in\mathbb{Z}$ to the latter inequality yields~\eqref{eqn:details_Delta}. We now evaluate the term $\Delta\boldsymbol{c}^{(\ell)}$ for an arbitrary $\ell$ based on the assumption that $\boldsymbol{c}^{(J)}$ is sampled from $f$ on the equispaced grid $2^{-J}\mathbb{Z}$. Since $f$ is bounded and differentiable, on the finest level, we have that
    \begin{equation}\label{eqn:Delta_c^J}
        \Delta\boldsymbol{c}^{(J)}\leq 2^{-J}\|f^\prime\|_{\infty}.
    \end{equation}
    This inequality is guaranteed by the Mean Value Theorem. In particular, for any $k\in\mathbb{Z}$ there exists a point $x_k$ in the open interval connecting the two consecutive grid points $k2^{-J}$ and $(k+1)2^{-J}$, such that $|c^{(J)}_{k+1}-c^{(J)}_{k}|\leq 2^{-J}|f^\prime (x_k)|$, and hence~\eqref{eqn:Delta_c^J} is obtained by taking the supremum over the index $k$. On coarser levels, however, since $\Delta$ commutes with the convolution operator, we get the recursive relation

    \begin{equation}\label{eqn:Delta_recursion}
        \Delta\boldsymbol{c}^{(\ell-1)} = \Delta\big(\boldsymbol{\gamma}^{(\ell)}*(\boldsymbol{c}^{(\ell)}\downarrow 2)\big)\leq \|\boldsymbol{\gamma}^{(\ell)}\|_1\cdot\Delta(\boldsymbol{c}^{(\ell)}\downarrow 2)\leq 2\|\boldsymbol{\gamma}^{(\ell)}\|_1\cdot\Delta\boldsymbol{c}^{(\ell)},
    \end{equation}
    which holds for $\ell=1,\dots,J$. Finally, iterative calculations of~\eqref{eqn:Delta_recursion} involving~\eqref{eqn:Delta_c^J} and~\eqref{eqn:details_Delta} yield~\eqref{details bound}.
\end{proof}

\end{document}
